\documentclass[11pt]{article}
\usepackage{franklin}
\renewcommand{\hdrtitle}{Final Exam --- Derivatives \& Antiderivatives}
\renewcommand{\hdrsub}{Sections 2.1--2.2 \& 5.1 \textbullet{} Fall 2026}
\renewcommand{\durationtext}{90 minutes}
\begin{document}
\makecover{Calculus \textbullet{} Final Examination}{Final Exam}{Derivatives \& Antidifferentiation}%
{Three Exercises --- Multiple Choice, Differentiation and Antidifferentiation\\ Scored out of 100 points}%
{Do not open this booklet until instructed. No notes, no calculators unless authorised.}

\section*{Instructions}
\studentbox
\bigskip
This examination covers the three course documents: \emph{Prerequisite Review \& Limits 1.5},
\emph{Derivative 2.1 \& Elementary Techniques of Differentiation 2.2}, and
\emph{Antidifferentiation 5.1}. It contains three exercises and is scored out of \textbf{100 points}.
\begin{itemize}[itemsep=2pt]
  \item \textbf{Exercise 1:} tick exactly \emph{one} box per question. No justification is required; a wrong or multiple answer scores 0.
  \item \textbf{Exercises 2--3:} show every step --- rewrite the function as powers of $x$, name the rule, then simplify. Unjustified answers receive no credit.
  \item Only the Constant, Power, Constant Multiple and Sum Rules are needed for derivatives: rewrite products, quotients and powers of binomials \emph{before} differentiating.
  \item For every indefinite integral include $+C$; for definite integrals give the exact value and a decimal to three places when it contains $e$ or $\ln$.
\end{itemize}
\scoretable

\newpage
\subsection*{Exercise 1 --- Multiple Choice \hfill {\normalsize\color{fnavy}30 points}}
\meta{(Level 3 --- Hard)\quad 5 questions $\times$ 6 points --- exactly one correct answer each}
\begin{enumerate}[label=\textbf{Q\arabic*.},leftmargin=*,itemsep=10pt]
\qcm{For $f(x)=\dfrac{1}{x-2}$, the simplified difference quotient $\dfrac{f(x+h)-f(x)}{h}$ is:}
  {$\dfrac{1}{(x-2)(x+h-2)}$}{$\dfrac{-h}{(x-2)(x+h-2)}$}{$\dfrac{-1}{(x-2)(x+h-2)}$}{$\dfrac{-1}{(x-2)^2}$}
\qcm{If $y=\dfrac{x^2-6\sqrt{x}}{x}$, then $y'(9)$ equals:}
  {$\tfrac89$}{$\tfrac{10}{9}$}{$0$}{$2$}
\qcm{Which function is an antiderivative of $f(x)=\dfrac{e^{3x}+2}{e^{x}}$?}
  {$\tfrac12 e^{2x}+2e^{-x}$}{$2e^{2x}-2e^{-x}$}{$\tfrac12 e^{2x}-2e^{-x}-5$}{$\dfrac{\frac13 e^{3x}+2x}{e^{x}}$}
\qcm{$\displaystyle\int_{4}^{1}\left(6\sqrt{x}-1\right)dx=$}
  {$25$}{$-31$}{$31$}{$-25$}
\qcm{Let $F'(x)=\dfrac{3}{x}-4x^3$ for $x>0$, with $F(1)=0$. Then $F(e)=$}
  {$3-e^4$}{$4+e^4$}{$-e^4$}{$4-e^4$}
\end{enumerate}

\newpage
\subsection*{Exercise 2 --- Differentiation \hfill {\normalsize\color{fnavy}40 points}}
\meta{(Level 2--3 --- Medium to Hard)\quad 5 questions $\times$ 8 points --- rewrite before you differentiate}
\begin{enumerate}[label=\textbf{(\alph*)},leftmargin=*,itemsep=4pt]
\item Find $f'(x)$ for $f(x)=2x^5-6\sqrt[4]{x^3}+\dfrac{4}{x^2}-e^3$. \pts{8}\answerspace{3.2cm}
\item Find $g'(x)$ for $g(x)=\left(\sqrt{x}-3\right)\left(x+\dfrac{2}{\sqrt{x}}\right)$. \pts{8}\answerspace{3.2cm}
\item Find $h'(x)$ for $h(x)=\dfrac{2x^3+5x-8}{x^2}$. \pts{8}\answerspace{3.2cm}
\item Find $k'(t)$ for $k(t)=(3t+2)^3$. \emph{(Expand first.)} \pts{8}\answerspace{3.2cm}
\item Let $f(x)=\dfrac{(x-3)^2}{\sqrt{x}}$. Find $f'(x)$, then the equation of the tangent line to the graph of $f$ at $x=9$. \pts{8}\answerspace{4cm}
\end{enumerate}

\newpage
\subsection*{Exercise 3 --- Antidifferentiation \hfill {\normalsize\color{fnavy}30 points}}
\meta{(Level 2 --- Medium)\quad 3 questions $\times$ 10 points}
\begin{enumerate}[label=\textbf{(\alph*)},leftmargin=*,itemsep=4pt]
\item Find $\displaystyle\int\left(9x^2+\frac{5}{x}-4e^{-2x}\right)dx$. \pts{10}\answerspace{5cm}
\item Find $\displaystyle\int\frac{(x-2)^2}{\sqrt{x}}\,dx$. \pts{10}\answerspace{5cm}
\item Evaluate $\displaystyle\int_{1}^{3}\left(3x^2+\frac{4}{x^2}-\frac{2}{x}\right)dx$. \pts{10}\answerspace{4cm}
\end{enumerate}
\end{document}
